\section{Introduction}

Protein structure prediction has long been recognized as one of the central challenges in computational biology, due to its critical role in understanding biological functions and developing novel therapeutics.~\cite{anfinsen1973principles} Traditional approaches, such as homology modeling and physics-based simulations, have achieved partial success but often suffer from limited accuracy or prohibitive computational cost. The release of AlphaFold2~\cite{jumper2021highly,evans2021protein} by DeepMind marked a major breakthrough in this field, delivering unprecedented accuracy in predicting three-dimensional protein structures from amino acid sequences. Its impact has extended across biology~\cite{perrakis2021ai}, protein design~\cite{goverde2023novo,wang2024protein}, and drug discovery~\cite{thornton2021alphafold,nussinov2023alphafold}, making large-scale deployment of AlphaFold2 inference increasingly desirable.
However, the adoption of AlphaFold2 in practical scenarios remains constrained by its high computational cost. A single inference typically requires substantial hardware resources, which limits its accessibility and large-scale deployment.

To address this challenge, recent efforts have focused on accelerating AlphaFold2 on GPUs. These works leverage techniques such as kernel fusion, memory optimization, and reduced-precision computation (e.g.,BF16) to improve inference efficiency. While these studies demonstrate substantial performance speedups, they are almost exclusively GPU-centric. In contrast, CPU-based acceleration of AlphaFold2 has received little attention. This lack of exploration is notable given that CPUs remain the dominant computing platform in many large-scale clusters, supercomputers, and cloud environments. Moreover, the implementations of AlphaFold2 are typically restricted to FP32 or BF16 precision, leaving aggressive quantization strategies such as FP16 and INT8 completely unexplored.

At the same time, modern CPUs are undergoing architectural evolution that brings them closer to accelerators in terms of matrix computation capability. For example, Intel has recently introduced the Advanced Matrix Extensions (AMX)~\cite{intel2022amx} in its latest processors, and Arm has developed Scalable Matrix Extensions (SME)~\cite{weidmann2021introducing} to enhance performance on modern many-core architectures. These hardware features enable CPUs to execute matrix-heavy operations with significantly higher throughput, narrowing the performance gap with GPUs for certain classes of workloads. Moreover, these matrix extensions are explicitly designed to support reduced-precision formats, making CPUs an emerging platform where quantization can play a central role in performance optimization.

The 72F processor offers a unique opportunity to revisit CPU-based acceleration of AlphaFold2. Unlike conventional CPUs, the 72F adopts a many-core architecture with eight NUMA domains. Each core is equipped with a dedicated matrix unit capable of high-throughput computation, while a high-bandwidth on-package memory sits between DDR and cache to mitigate the memory bottleneck. 
Combined with its native support for FP16 and INT8 computation, the 72F is particularly well-suited for exploring the impact of quantization on large-scale models such as AlphaFold2.

% 矩阵单元一般对与低精度计算有更好的支持。在Intel AMX中，单条BF16矩阵指令可以执行1024次运算，而INT8则可以执行2048次运算。在ARM SME中，单条INT8外积指令也可以执行FP16或BF16两倍的计算量。因此量化是提升矩阵单元利用率的关键技术。
% Model quantization has emerged as a key technique to improve the efficiency of deep learning inference by reducing both memory footprint and computational cost. By converting high-precision floating-point operations into low-precision integer arithmetic, quantization enables better utilization of hardware resources such as matrix units, while also reducing memory bandwidth pressure. Recent studies have demonstrated that INT8 quantization can achieve substantial speedups with negligible accuracy degradation for vision and language models.

% Despite its effectiveness, quantization has not been systematically applied to AlphaFold2. The model's heterogeneous computation patterns, including attention, triangular updates, and geometry-based transformations, make straightforward quantization non-trivial. Moreover, maintaining structural accuracy under low-precision computation is particularly challenging for AlphaFold2, as minor numerical errors can accumulate and lead to significant deviations in predicted atomic coordinates. These challenges highlight the need for specialized quantization strategies that are tailored to AlphaFold2's architecture and workloads, particularly in the context of CPU execution with matrix units.

% To address these challenges, we introduce MatrixFold, an efficient implementation of AlphaFold2 for inference. MatrixFold is designed to fully exploit the capabilities of modern many-core CPUs with matrix units, while incorporating INT8 quantization to maximize throughput. Our contributions are as follows:
Optimizing AlphaFold2 on the 72F CPU presents several challenges. First, the model is dominated by large-scale dense linear algebra, requiring efficient matrix multiplication kernels that can fully exploit the matrix units of each core. Second, the processor's many-core, multi-NUMA design necessitates careful handling of communication across memory domains to avoid bandwidth contention and excessive synchronization. Third, the adoption of reduced-precision inference on CPUs introduces new complexity, as FP16 and INT8 quantization must be integrated without sacrificing model accuracy. Finally, the variable length of protein sequences makes parallelization non-trivial: overly aggressive parallelism on short sequences can lead to small per-thread workloads, synchronization overheads, and costly cross-NUMA communication.

To overcome these challenges, we make several design choices. We develop a high-performance GEMM implementation tailored to the 72F's matrix units, ensuring that the dominant computation is efficiently executed. We design a shared-memory communication mechanism optimized for multi-NUMA systems, reducing inter-node data transfer costs. We enable FP16 inference and implement INT8 quantization, providing the first demonstration of quantized AlphaFold2 inference on CPUs. Furthermore, we introduce a parallelism detection mechanism that adaptively selects thread counts to balance workload granularity against synchronization and communication overhead.
In this work, we present MatrixFold, a systematic study of AlphaFold2 optimization on CPUs, targeting the 72F architecture. Our contributions are as follows:
\begin{itemize}
    \item We present a systematic study and implementation of mix-precision (FP16 and INT8) AlphaFold2 inference on a many-core CPU architecture. 
    \item We propose a suite of computation and communication optimizations specifically tailored for the target 72F processor, leveraging its dedicated matrix units and multi-NUMA architecture.
    \item We conduct a comprehensive evaluation on the CASP14 dataset, demonstrating that MatrixFold achieves a significant performance advantage on long sequences, outperforming the state-of-the-art GPU implementation, FastFold on an NVIDIA A800 GPU, by 1.2$\times$ to 3.2$\times$. 
    \item We validate that our mixed-precision  strategy achieves substantial performance gains without compromising prediction accuracy, confirming the viability and effectiveness of low-precision inference for AlphaFold2 on CPUs.
    % \item We design a set of high-performance GEMM kernels leveraging matrix units to accelerate the inference of AlphaFold2. 
    % \item We enable quantized inference on CPUs by implementing both FP16 and INT8 execution, moving beyond the conventional FP32/BF16 precision levels.
    % \item We design NUMA-aware parallelization and memory optimizations that exploit the 72F's many-core design and high-bandwidth on-package memory.
    % \item We provide a comprehensive evaluation demonstrating that our approach achieves substantial performance improvements while maintaining accuracy.
\end{itemize}